\section{The Water Melting Point as a Proton Transfer Eigenmode}
\label{sec:melting_eigenmode}


\begin{objectivebox}
\begin{itemize}
    \item Derive the water melting point ($T_m = 273.46$~K, $+0.14\%$ error) as the anharmonic saturation eigenmode of the proton-transfer $LC$ oscillator.
    \item Understand how all seven bulk water properties emerge from a single proton $LC$ eigenmode with zero free parameters.
    \item Explain the density anomaly as a competition between melting contraction and thermal expansion near the theoretical melting eigenmode.
\end{itemize}
\end{objectivebox}

\subsection{The Physical Problem}

Water's melting point ($T_m = 273.15$~K) and the anomalous density maximum 
at approximately $4^{\circ}$C are macroscopic observables that arise from the 
microscopic structure of the hydrogen-bond network. Classical thermodynamics 
treats $T_m$ as an empirical constant; the AVE framework derives it as the 
\emph{eigenfrequency} of the proton cavity formed by the O--H$\cdots$O bridge.

\subsection{Step 1: LC Analogs of the H-Bond Network}

Water's tetrahedral network ($z = 4$) is modelled as a discrete LC lattice 
where each node (H$_2$O molecule) connects to four neighbours via hydrogen bonds.

\begin{center}
\small
\renewcommand{\arraystretch}{1.2}
\begin{tabularx}{\textwidth}{@{}lccX@{}}
\toprule
\textbf{Element} & \textbf{$\mu$ (Inertia)} & \textbf{$\varepsilon$ (Compliance)} & \textbf{$Z = \sqrt{\mu/\varepsilon}$} \\
\midrule
O--H covalent bond & $\mu_\text{OH}$ (reduced) & $1/k_\text{OH}$ & Intramolecular $Z$ \\
H-bond             & $m_{\text{H}_2\text{O}}$ & $1/k_\text{hb}$ & Intermolecular $Z$ \\
Lattice ($z = 4$)  & $m_{\text{H}_2\text{O}}$ (loaded) & $1/(z \cdot k_\text{hb})$ & Network $Z_0$ \\
\bottomrule
\end{tabularx}
\end{center}

\subsection{Step 2: H-Bond Spring Constant from Op4}

The H-bond energy $E_\text{hb}$ and equilibrium distance $d_\text{hb}$ are derived 
from the universal pairwise potential (Op4). The coupling constant for the H-bond 
is set by the impedance mismatch (Op3) between oxygen and hydrogen port radii:
\begin{equation}
  K_\text{hb} = \Gamma^2 \alpha \hbar c, 
  \qquad \Gamma = \frac{r_\text{O} - r_\text{H}}{r_\text{O} + r_\text{H}}
  \label{eq:hb_coupling}
\end{equation}
where $r_\text{O}$ and $r_\text{H}$ are the atomic port impedances (Op1/Op3).

The H-bond equilibrium distance $d_\text{hb}$ is the minimum of $U(r)$ from Op4, 
and the harmonic spring constant is:
\begin{equation}
  k_\text{hb} = \frac{2 E_\text{hb}}{d_\text{hb}^2}
  \label{eq:k_hb}
\end{equation}

\noindent\textbf{Dimensional check:}
\begin{equation}
  [k_\text{hb}] = \frac{[\text{J}]}{[\text{m}^2]} 
  = \frac{[\text{kg}\cdot\text{m}^2/\text{s}^2]}{[\text{m}^2]} 
  = [\text{kg}/\text{s}^2] = [\text{N/m}] \quad\checkmark
\end{equation}

\noindent\textbf{Numerical evaluation:}
\begin{equation}
  k_\text{hb} = \frac{2 \times 3.458 \times 10^{-20}}{(1.754 \times 10^{-10})^2} 
  = 2.248 \text{ N/m}
\end{equation}

\subsection{Step 3: O--H Spring Constant from the Coulomb Bond Solver}

The intramolecular O--H spring constant $k_\text{OH}$ is derived from the 
Coulomb bond force constant solver, which computes $d^2E/dr^2$ at equilibrium 
for the covalent bond using Slater orbitals and three lattice topology corrections:

\begin{enumerate}
  \item \textbf{Isotropy projection} ($1/3$): Bond stretching acts along 
        one of three equivalent spatial dimensions on the isotropic SRS lattice 
        (Axiom~1).
  \item \textbf{Three-phase balance} ($1/\sqrt{3}$ per terminal atom): 
        Each interior lattice node is a three-connected Wye junction. 
        Terminal atoms (H) are unbalanced single-phase loads.
  \item \textbf{Lone-pair dynamic softening}: The four lone-pair electrons on oxygen 
        couple weakly ($\cos^2(109.5^{\circ}) = 1/9$) to the stretching mode, 
        broadening the potential well.
\end{enumerate}

The combined correction factor for the O--H bond (one terminal H) is:
\begin{equation}
  f_\text{corr} = \frac{1}{3} \times \frac{1}{\sqrt{3}} 
  = \frac{1}{3\sqrt{3}} \approx 0.1925
\end{equation}

Applied to the raw Coulomb curvature:
\begin{equation}
  k_\text{OH} = f_\text{corr} \times 
  \left.\frac{d^2 E}{d r^2}\right|_{r = r_\text{eq}} 
  = 791 \text{ N/m}
  \label{eq:k_OH}
\end{equation}

The experimental value is $745 \pm 5$~N/m (from the O--H symmetric stretch at 
$\tilde{\nu} = 3657$~cm$^{-1}$), giving an engine error of $+6.2\%$.

\subsection{Step 4: The Proton Transfer Eigenmode}

The melting transition is identified with the fundamental eigenmode of the 
O--H$\cdots$O bridge---the hydrogen atom oscillating between two oxygen atoms 
through a series spring system:

\begin{equation}
  \underbrace{\text{O}_1}_{\text{covalent}} 
  \overset{k_\text{OH}}{=\!=\!=} 
  \underbrace{\text{H}}_{{m_\text{H}}} 
  \overset{k_\text{hb}}{\cdots\cdots\cdots} 
  \underbrace{\text{O}_2}_{\text{H-bonded}}
  \label{eq:bridge_diagram}
\end{equation}

The two springs act in \emph{series} for the proton displacement:
\begin{equation}
  k_\text{series} = \frac{k_\text{OH} \cdot k_\text{hb}}{k_\text{OH} + k_\text{hb}}
  \label{eq:k_series}
\end{equation}

\noindent\textbf{Dimensional check:}
\begin{equation}
  [k_\text{series}] = \frac{[\text{N/m}] \cdot [\text{N/m}]}{[\text{N/m}] + [\text{N/m}]} 
  = \frac{[\text{N}^2/\text{m}^2]}{[\text{N/m}]} 
  = [\text{N/m}] \quad\checkmark
\end{equation}

Since $k_\text{OH} = 791$~N/m $\gg k_\text{hb} = 2.248$~N/m:
\begin{equation}
  k_\text{series} = \frac{791 \times 2.248}{791 + 2.248} = 2.242 \text{ N/m} \approx k_\text{hb}
\end{equation}

The soft H-bond spring completely dominates the series combination. This is a 
general result: the weakest link sets the eigenfrequency.

The eigenfrequency of the proton oscillating in this cavity is:
\begin{equation}
  \omega_m = \sqrt{\frac{k_\text{series}}{m_\text{H}}}
  \label{eq:omega_m}
\end{equation}

\noindent\textbf{Dimensional check:}
\begin{equation}
  [\omega_m] = \sqrt{\frac{[\text{N/m}]}{[\text{kg}]}} 
  = \sqrt{\frac{[\text{kg}/\text{s}^2]}{[\text{kg}]}} 
  = \sqrt{[\text{s}^{-2}]} = [\text{rad/s}] \quad\checkmark
\end{equation}

The melting condition is $kT_m = \hbar\omega_m$---the thermal energy per mode 
equals one quantum of the proton cavity oscillation. Below this temperature, 
the proton cannot classically hop between donor and acceptor oxygen sites, 
and the H-bond network retains directional coherence (solid phase):
\begin{equation}
  \boxed{T_m = \frac{\hbar}{k_B} \sqrt{\frac{k_\text{series}}{m_\text{H}}}}
  \label{eq:Tm_eigenmode}
\end{equation}

\noindent\textbf{Dimensional check:}
\begin{equation}
  [T_m] = \frac{[\text{J}\cdot\text{s}]}{[\text{J/K}]} \cdot [\text{s}^{-1}]
  = \frac{[\text{kg}\cdot\text{m}^2/\text{s}]}{[\text{kg}\cdot\text{m}^2/(\text{s}^2\cdot\text{K})]}
  \cdot [\text{s}^{-1}] = [\text{K}] \quad\checkmark
\end{equation}

\subsection{Step 5: Numerical Evaluation (Bare Eigenmode)}

\begin{center}
\small
\renewcommand{\arraystretch}{1.2}
\begin{tabularx}{\textwidth}{@{}lllX@{}}
\toprule
\textbf{Quantity} & \textbf{Value} & \textbf{Source} & \textbf{Type} \\
\midrule
$E_\text{hb}$ & $3.458 \times 10^{-20}$ J & Op4 well depth & Engine-derived \\
$d_\text{hb}$ & 1.754 \AA & Op4 equilibrium & Engine-derived \\
$k_\text{hb}$ & 2.248 N/m & $2E_\text{hb}/d_\text{hb}^2$ & Engine-derived \\
$k_\text{OH}$ & 791 N/m & Coulomb bond solver & Engine-derived \\
$k_\text{series}$ & 2.242 N/m & Eq.~\eqref{eq:k_series} & Engine-derived \\
$m_\text{H}$ & $1.674 \times 10^{-27}$ kg & Proton mass & Fundamental \\
$\hbar$ & $1.055 \times 10^{-34}$ J$\cdot$s & & Fundamental \\
$k_B$ & $1.381 \times 10^{-23}$ J/K & & Fundamental \\
\midrule
$\omega_0$ & $3.659 \times 10^{13}$ rad/s & Eq.~\eqref{eq:omega_m} & Derived \\
$T_0$ (bare) & \textbf{279.5 K ($+6.3^{\circ}$C)} & Eq.~\eqref{eq:Tm_eigenmode} & Derived \\
$T_m$ (experiment) & 273.15 K ($0^{\circ}$C) & & \\
\textbf{Bare error} & $\mathbf{+2.3\%}$ & & \\
\bottomrule
\end{tabularx}
\end{center}

This bare eigenmode $T_0 = 279.5$~K is the \emph{harmonic} LC resonance. It accurately predicts the temperature of maximum density ($T(\rho_\text{max}) = T_0(1-\alpha) = 4.26^{\circ}$C, experiment: $3.98^{\circ}$C), but overestimates the melting point by $+2.3\%$ because it neglects the Op4 anharmonicity.

\subsection{Step 6: Axiom~4 Saturation Correction}

The $+2.3\%$ residual is resolved by Axiom~4. At finite temperature, the proton oscillation amplitude $\delta x = \sqrt{k_B T / k_\text{hb}}$ is a significant fraction of the well width $d_\text{hb}$. The thermal strain is:
\begin{equation}
  r = \sqrt{\frac{k_B T}{2\,E_\text{hb}}}
\end{equation}
At $T_0 = 279.5$~K: $r = 0.236$, which is well past the Regime~I/II boundary at $\sqrt{2\alpha} = 0.121$. The Op4 potential is \emph{not} perfectly parabolic in Regime~II; Axiom~4 provides the saturation function:
\begin{equation}
  S(r) = \sqrt{1 - r^2}
\end{equation}
This softens the spring constant: $k_\text{hb}^\text{eff}(T) = k_\text{hb}\,S(r(T))$. The melting point is the \textbf{self-consistent} solution:
\begin{equation}
  T_\text{sat} = \frac{\hbar}{k_B}\sqrt{\frac{k_\text{hb}^\text{eff}(T_\text{sat})}{m_\text{H}}}
\end{equation}
Iterating from $T_0$ converges in $\sim$13 steps to $T_\text{sat} = 275.5$~K ($r = 0.234$, $S = 0.972$).

Applying the Axiom~2 vacuum--lattice loading $(1-\alpha)$:
\begin{equation}
\boxed{
  T_\text{melt} = T_\text{sat}(1-\alpha) = 275.5 \times 0.99270 = 273.5 \text{ K} \quad (+0.14\%)
}
\end{equation}

The correction chain is:
\begin{center}
\small
\renewcommand{\arraystretch}{1.2}
\begin{tabularx}{\textwidth}{@{}lcX@{}}
\toprule
\textbf{Stage} & \textbf{$T_m$ (K)} & \textbf{Error} \\
\midrule
Bare eigenmode $T_0$ & 279.5 & $+2.31\%$ \\
$+$ Axiom~4 saturation & 275.5 & $+0.88\%$ \\
$+$ Axiom~2 $\alpha$-loading & 273.5 & $+0.14\%$ \\
\bottomrule
\end{tabularx}
\end{center}

\subsection{Lattice Loading Analysis}

A natural question arises: does the local lattice density modify the eigenfrequency?

Each oxygen atom in the tetrahedral lattice is held by $(z - 1) = 3$ additional 
H-bonds. A coupled 2-mass model (O$_1$+H covalent unit vs.\ O$_2$ anchored by 
$3 k_\text{hb}$ lattice springs) was solved as a generalized matrix eigenvalue 
problem:
\begin{equation}
  \mathbf{K} \mathbf{X} = \omega^2 \mathbf{M} \mathbf{X}
\end{equation}
with stiffness matrix $\mathbf{K}$ having entries $[z k_\text{hb}, -k_\text{hb}; 
-k_\text{hb}, z k_\text{hb}]$ and mass matrix $\mathbf{M} = 
\text{diag}(m_\text{O} + m_\text{H},\ m_\text{O})$.

The result: the lattice optical mode appears at $T \approx 155$~K---this is the 
\emph{O--O stretching} frequency, not the proton transfer mode. The proton 
transfer mode does not appear because $k_\text{OH} \gg k_\text{hb}$ locks the 
proton rigidly to O$_1$.

The lattice correction to $T_m$ is bounded by the mass ratio:
\begin{equation}
  \frac{\Delta T_m}{T_m} \sim \frac{m_\text{H}}{2 m_\text{O}} 
  = \frac{1.008}{2 \times 15.999} = 0.032 \approx 3\%
\end{equation}
Applied to the $6.3$~K error, this gives $\Delta T_m \lesssim 0.2$~K---negligible. 
The O atoms are effectively clamped by their mass ratio alone; the lattice springs 
make them \emph{more} clamped, not less.

\subsection{Physical Interpretation}

The lattice melts when the thermal energy per mode ($kT$) reaches the energy 
quantum of the proton oscillating in the O--H$\cdots$O bridge ($\hbar\omega_m$). 
At this temperature, the proton can classically hop between donor and acceptor 
oxygen sites, destroying the directional coherence of the H-bond network.

This is isomorphic to atomic ionization: ionization energies are eigenvalues of 
the \emph{electron cavity}; the melting point is the eigenvalue of the 
\emph{proton cavity}. The mechanism is identical across scales---the only 
difference is the mass of the oscillating particle and the stiffness of the 
confining well.

\subsection{Residual Error Analysis}

The $+2.3\%$ bare eigenmode residual is fully resolved by two corrections, both derived from the AVE axioms:
\begin{enumerate}
  \item \textbf{Axiom~4 saturation} ($+2.31\% \to +0.88\%$): the Op4 spring softens at the thermal strain $r = 0.234$ (Regime~II). Self-consistent solution gives $T_\text{sat} = 275.5$~K.
  \item \textbf{Axiom~2 $\alpha$-loading} ($+0.88\% \to +0.14\%$): the vacuum lattice impedance loads the eigenmode by $(1-\alpha)$.
\end{enumerate}

The final residual of $+0.14\%$ ($0.38$~K) is within the uncertainty of the Op4 well parameters and the parabolic approximation to the well curvature.

\begin{resultbox}{Corrected Melting Point}
\begin{equation}
  T_\text{melt} = T_\text{sat}(1-\alpha) = 273.5 \text{ K} = 0.38^{\circ}\text{C}
  \qquad (\text{exp: } 273.15 \text{ K}, \; \text{error: } +0.14\%)
\end{equation}
\end{resultbox}

\subsection{Seven Observables from One Spring Constant}

The proton eigenmode derived in this section is the foundation for a broader result: \emph{seven} independent macroscopic observables of bulk water are derivable from the single input $k_\text{hb} = 2.247$~N/m. The complete proof---including speed of sound ($+0.5\%$), temperature of maximum density ($+0.28^{\circ}$C), enthalpy of vaporisation ($+1.7\%$), surface tension ($-1.1\%$), boiling point ($-0.15\%$), and specific heat ($-0.7\%$)---is presented in Volume~III, Chapter~\ref{ch:water_lc_lattice}.

\begin{summarybox}
\begin{itemize}
  \item The water melting point is derived as the self-consistent saturated proton 
        eigenmode: $T_\text{melt} = T_\text{sat}(1-\alpha) = 273.5$~K ($+0.14\%$).
  \item The bare eigenmode $T_0 = 279.5$~K sets the temperature of maximum density
        via $T(\rho_\text{max}) = T_0(1-\alpha) = 4.26^{\circ}$C ($+0.28^{\circ}$C).
  \item The H--O--H bond angle is derived from the sp$^3$ tetrahedral angle 
        via Op3 impedance transmission: $\theta = \arccos(-1/(3(1+\Gamma))) = 104.48^{\circ}$
        ($+0.03^{\circ}$ from experiment).
  \item The Axiom~4 saturation correction $S(r) = \sqrt{1-r^2}$ resolves the 
        $+2.3\%$ bare residual by accounting for spring softening in Regime~II.
  \item All inputs are engine-derived or fundamental constants; zero spectroscopic 
        or empirical inputs are required.
  \item The same spring constant predicts seven independent observables 
        (see Volume~III, Chapter~\ref{ch:water_lc_lattice}).
\end{itemize}
\end{summarybox}

\begin{exercisebox}
\begin{enumerate}
  \item Verify that $k_\text{series} \approx k_\text{hb}$ by expanding 
        Eq.~\eqref{eq:k_series} to first order in the small parameter 
        $k_\text{hb}/k_\text{OH}$.
  \item Compute $T_m$ for a hypothetical D$_2$O (heavy water) lattice by replacing 
        $m_\text{H}$ with $m_\text{D} = 2.014$~amu. Compare with the experimental 
        value $T_m(\text{D}_2\text{O}) = 276.97$~K.
  \item Derive the lattice loading correction to $T_m$ from the 2-mass eigenvalue 
        problem. Show that the correction scales as $m_\text{H}/(2m_\text{O})$.
  \item Compute the H--S--H bond angle for H$_2$S using $\Gamma_\text{SH} = 
        (r_\text{S} - r_\text{H})/(r_\text{S} + r_\text{H})$ and compare with the 
        experimental value of $92.1^{\circ}$.
  \item Show that the small-signal (Op3) and large-signal (Op8) bond angle 
        formulas converge to within $0.19^{\circ}$ for H$_2$O, and explain 
        physically why this convergence implies a self-consistent eigenvalue.
\end{enumerate}
\end{exercisebox}

%% == Section: Bond Angle Derivation =====================================

\section{The H--O--H Bond Angle as an Impedance Eigenvalue}
\label{sec:bond_angle_derivation}

The preceding section treated the melting point as the eigenfrequency of the 
O--H$\cdots$O bridge. A second geometric constant---the intramolecular 
H--O--H bond angle $\theta_\text{HOH} = 104.45^{\circ}$---was until now accepted 
as an empirical input. This section derives it from the same axioms.

\subsection{The sp$^3$ Tetrahedral Angle (Axiom~1)}

The K4 unit cell of the SRS lattice (Axiom~1) has four-connected nodes. 
The equilibrium angle between any two bonds at a four-connected node is the 
tetrahedral angle:
\begin{equation}
  \cos\theta_\text{tet} = -\frac{1}{3}, \qquad 
  \theta_\text{tet} = \arccos\!\left(-\tfrac{1}{3}\right) = 109.47^{\circ}
  \label{eq:theta_tet}
\end{equation}
This assumes all four orbitals are equivalent (equal impedance, equal angular 
weight). Oxygen's valence shell has four sp$^3$ orbitals, but they are 
\emph{not} equivalent: two are bonding (loaded by H) and two are lone pairs 
(unloaded).

\subsection{Op3 Small-Signal Correction}

The O--H junction has a reflection coefficient (Op3):
\begin{equation}
  \Gamma = \frac{r_\text{O} - r_\text{H}}{r_\text{O} + r_\text{H}} = \frac{1}{3}
  \label{eq:Gamma_OH}
\end{equation}
where $r_\text{O}$ and $r_\text{H}$ are the atomic port impedances.

The bonding orbital is a \emph{partially transmitted} wave through the O--H 
junction: the proton loads the orbital, and the transmission factor $(1 + \Gamma)$ 
determines the effective angular weight of the bonding pair relative to the 
lone pair (which is a totally reflected wave, $\Gamma = -1$, with weight $1$).

The cosine of the bond angle is the tetrahedral cosine divided by 
the transmission factor:
\begin{resultbox}{H--O--H Bond Angle (Op3 Small-Signal)}
\begin{equation}
  \cos\theta_\text{HOH} = \frac{\cos\theta_\text{tet}}{1 + \Gamma} 
  = \frac{-1/3}{1 + 1/3} = -\frac{1}{4}
  \label{eq:theta_HOH}
\end{equation}
\begin{equation}
  \theta_\text{HOH} = \arccos\!\left(-\tfrac{1}{4}\right) = 104.48^{\circ}
  \qquad (\text{exp: } 104.45^{\circ}, \;\text{error: } +0.03^{\circ})
\end{equation}
\end{resultbox}

\noindent\textbf{Dimensional check:} All quantities are dimensionless (ratios 
of port impedances). The cosine is bounded $|\cos\theta| \leq 1$. 
For any $\Gamma \in [0, 1)$: $|\cos\theta| = 1/(3(1+\Gamma)) < 1/3 < 1\ \checkmark$.

\subsection{Op8 Large-Signal Confirmation}

The FCC packing fraction $\varphi = \pi\sqrt{2}/6 \approx 0.7405$ (Op8/Axiom~2) 
provides an independent large-signal check. The lone pairs fill the packed 
volume $\varphi$ of the sp$^3$ angular space, compressing the bonding angle by 
the same fraction:
\begin{equation}
  \cos\theta_\text{HOH}^{(\text{large})} = \cos\theta_\text{tet} \times \varphi
  = -\frac{\varphi}{3} = -\frac{\pi\sqrt{2}}{18} \approx -0.2468
  \label{eq:theta_HOH_large}
\end{equation}
\begin{equation}
  \theta_\text{HOH}^{(\text{large})} = \arccos\!\left(-\tfrac{\pi\sqrt{2}}{18}\right) 
  = 104.29^{\circ} \qquad (\text{exp: } 104.45^{\circ}, \;\text{error: } -0.16^{\circ})
\end{equation}

The two estimates\,---\,one from local junction impedance (Op3), one from 
bulk packing geometry (Op8)\,---\,converge to within $0.19^{\circ}$. This 
convergence between small-signal and large-signal analyses is the hallmark of 
a self-consistent eigenvalue: the molecule's local impedance structure is 
compatible with the global lattice packing constraint.

\begin{center}
\small
\renewcommand{\arraystretch}{1.2}
\begin{tabularx}{\textwidth}{@{}lccX@{}}
\toprule
\textbf{Method} & \textbf{Formula} & \textbf{$\theta$ (deg)} & \textbf{Error} \\
\midrule
Bare sp$^3$ (Axiom~1) & $\arccos(-1/3)$ & 109.47 & $+5.02^{\circ}$ \\
Op3 small-signal & $\arccos(-1/(3(1+\Gamma)))$ & 104.48 & $+0.03^{\circ}$ \\
Op8 large-signal & $\arccos(-\varphi/3)$ & 104.29 & $-0.16^{\circ}$ \\
Experimental & & 104.45 & --- \\
\bottomrule
\end{tabularx}
\end{center}

\subsection{Physical Interpretation}

The deviation from the tetrahedral angle has an exact EE analog. A four-port 
junction (sp$^3$ node) with two loaded ports (bonding pairs, terminated in 
$Z_\text{H}$) and two open ports (lone pairs, $\Gamma = -1$) has an asymmetric 
scattering matrix. The loaded ports have lower input impedance, which allows the 
unloaded ports to widen their angular separation. The bonding pairs are squeezed.

This is structurally identical to the impedance pulling of a loaded resonator: 
attaching a load to a cavity shifts its resonant frequency. Here, attaching 
hydrogen atoms to two of four sp$^3$ orbitals shifts their angular separation 
from $\arccos(-1/3) \to \arccos(-1/4)$.

Crucially, $\Gamma = 1/3$ for the O--H junction makes the formula exact:
\begin{equation}
  \cos\theta = -\frac{1}{3} \to -\frac{1}{3 \times 4/3} = -\frac{1}{4}
\end{equation}
The angle transitions from $\arccos(-1/N)$ with $N = 3$ (tetrahedral) to 
$N = 4$ (loaded tetrahedral). The loaded tetrahedral angle is the next rational 
cosine in the sequence $\arccos(-1/N)$.

%% == Section: Topological Cell Collapse ==================================

\section{Topological Cell Collapse (State II Volume)}
\label{sec:topological_cell_collapse}

The classic thermodynamic anomaly of water---that it shrinks upon melting to reach a maximum density at $\sim 4^\circ$C---arises directly from the yielding of its topological LC matrix. 

Classical models struggle to reproduce the density of liquid water (998~kg/m$^3$ at 20$^\circ$C) from its molecular geometry without invoking arbitrary empirical packing radii. If one treats a water molecule as a continuous solid boundary (like a 3-sphere exact geometric union), its bare spatial volume is rigidly $5.56\ \AA^3$. Dividing this by the $\varphi_{FCC}$ Kepler packing fraction yields an absurd packed density of nearly 4000~kg/m$^3$.

In AVE (Axiom~1), spatial volume is not the geometric hull of a particle, but the volumetric capacity of a loaded discrete network cell. We derive the macroscopic liquid state (State II) volume without spatial bounds, relying strictly on the K=2G topological collapse (Axiom~2).

\subsection{State I: The Open Lattice ($V_I$)}
The rigid, unyielded State I is a tetrahedral lattice dictated by the Op4 pairwise minimum distance ($d_{OO} = d_{OH} + d_{hb} = 2.726\ \AA$). The unit cell volume natively occupied by one molecule in this expanded topology is:
\begin{equation}
    V_I = \frac{8}{3\sqrt{3}} d_{OO}^3 = 31.21\ \AA^3
\end{equation}
This establishes a theoretical ice density of 958~kg/m$^3$ (experimental Ice Ih is 917~kg/m$^3$).

\subsection{State II: FCC Maximal Yield ($V_{II}$)}
When the thermal energy surpasses the Op4 barrier ($f_I < 1.0$), the strict tetrahedral orientation fractures. The network yields to the universal packing ceiling prescribed by Axiom~2: $K=2G \rightarrow$ FCC packing.

Crucially, the foundational node-to-node proximity ($d_{OO} = 2.726\ \AA$) does not change---the Op4 minimum remains intact even in the disorganized slush. The only metric that changes is the global topology: the exact spatial cells of State I forfeit their stretched spatial arrangement and collapse into FCC maximal packing. Therefore, the volume of the fully yielded State II fluid is simply the State I volume compacted by Kepler's packing fraction $\varphi = \pi\sqrt{2}/6 \approx 0.7405$:
\begin{resultbox}{Topological Cell Collapse (State II Volume)}
\begin{equation}
    V_{II} = V_I \times \varphi_{FCC} = V_I \times \frac{\pi\sqrt{2}}{6} \approx 23.11\ \AA^3
\end{equation}
\end{resultbox}

\subsection{The Density Anomaly}
By evaluating the structural mixture at $20^\circ$C, the lattice survival fraction (from the Op4 cooperative solver) is $f_I \approx 0.735$. Incorporating anharmonic thermal expansion ($V_{th} \approx 0.67\ \AA^3$):
\begin{equation}
    V_{avg} = f_I V_I + (1-f_I) V_{II} + V_{th} 
    \approx 0.735(31.21) + 0.265(23.11) + 0.67 = 29.74\ \AA^3
\end{equation}
The resulting macroscopic density is $m_{H_2O} / V_{avg} = 1005.8$~kg/m$^3$. This yields precisely the experimental liquid state with a $\mathbf{+0.7\%}$ error, entirely devoid of empirical collision or hard-sphere parameters.

The density anomaly occurs because the melting term $(V_I \to V_{II})$ is a pure contraction, while thermal vibrations $(V_{th})$ are pure expansions. The competition balances precisely near the theoretical melting eigenmode.

